% New working identifiers A.DC1--A.DC4, allocated paper numbers A.45--A.48.
% GPT-6 Astra Pro, Pro mode; separate effort setting not exposed.
% This is a direct algebraic replacement for the characteristic-zero route.
\subsection{A direct algebraic configuration argument}

We give a proof in characteristic $p$ whose inputs contain neither lifting
of curves to characteristic zero nor the specialization statements of
SGA~1, Expos\'e XIII, 2.9--2.11. It also allows $p\in\Sigma$.
We retain the definitions of $F_{\mathrm{cr}}$ and $E_{\mathrm{cr}}$ in
A.CF2, now using geometric \emph{\'etale} fundamental groups, with
$r\in\Sigma$ and $r\ne p$.

\begin{claim}{A.DC1}{Proved}{Finite single-prime torsors in a constant compactification}
Let $k$ be algebraically closed of characteristic $p\geq0$, let $r$ be
a prime invertible in $k$, and let $C_0,S$ be smooth proper connected
$k$-curves, with $g(S)\geq1$. Let $m\geq1$, and let
$d_1,\ldots,d_m:S\to C_0\times S$ be pairwise disjoint sections.
Put $X=C_0\times S$, $D=\coprod_i d_i(S)$, $U=X-D$, and
$f:U\to S$. For a geometric point $s$ of $S$, every finite $r$-group
quotient of $\pi_1(U_s)$ is dominated by the restriction of a finite
$r$-group quotient of $\pi_1(U)$. Equivalently,
$\pi_1(U_s)^{(r)}\to\pi_1(U)^{(r)}$ is injective.
More precisely, every finite $r$-group torsor on $U_s$ extends to a
torsor with the same group on $U\times_S S'$ for some connected
finite \'etale cover $S'\to S$ whose Galois closure is an $r$-cover.
\end{claim}

\begin{proof}
Write $a=\mathbf F_r$ for the constant sheaf and fix an identification
$a(1)\simeq a$, possible over $k$. All cohomology in this proof is
\'etale cohomology unless continuous group cohomology is explicitly
indicated. We use the following algebraic inputs, in the exact forms
recorded in our source table: cohomological purity for a smooth pair of
codimension one (\citealp[Theorem 16.7]{CF_Milne2013}), proper base change
(ibid., Theorem 17.10), the cohomology of a proper smooth curve (ibid.,
Proposition 14.2), Leray (ibid., Theorem 12.7), and Cartan--Leray for
finite \'etale covers and their inverse limits (ibid., Theorem 14.9 and
Remark 14.10). None of these inputs is a characteristic-zero comparison.

\paragraph{The relative abelian calculation.}
Let $j:U\hookrightarrow X$ and $i:D\hookrightarrow X$.
Purity and localization give the distinguished triangle
\[
 i_*a_D(-1)[-2]\longrightarrow a_X\longrightarrow Rj_*a_U
 \longrightarrow i_*a_D(-1)[-1].
\]
Push it forward by the proper projection $\bar f:X\to S$.
Proper base change identifies $R^q\bar f_*a$ with the constant sheaf
$H^q(C_0,a)_S$. Since $D$ is a disjoint union of $m$ sections,
its higher direct images vanish in positive degrees. Thus
\[
 R^0f_*a=a_S,\qquad R^qf_*a=0\ (q\geq2),
\]
and, with $\mathcal M=R^1f_*a$,
\[
 0\to H^1(C_0,a)_S\to\mathcal M
 \to a(-1)_S^{\,m}\overset{\sum}{\longrightarrow}a(-1)_S\to0.       \tag{DC.1}
\]
The last map sends every labelled point class to its degree, namely one.
It is surjective even when $r\mid m$. In particular $\mathcal M$ is
a finite locally constant sheaf. The two terms in the resulting short
exact sequence
\[
 0\to H^1(C_0,a)_S\to\mathcal M
 \to\ker(a^m\overset{\sum}{\longrightarrow}a)_S(-1)\to0
\]
are constant. After a vector-space splitting at one geometric point,
monodromy is a subgroup of the additive group of homomorphisms from the
last vector space to the first. Its image is therefore an elementary
abelian $r$-group. No equivariant splitting and no division by $m$ is
used.

Here the identifications with the cohomology of $U_s$ are themselves
part of the argument, not an unproved base-change assertion for an open
morphism. Pulling the displayed localization triangle to a geometric
fibre gives its localization triangle: a section is transverse to the
fibre with intersection multiplicity one, so its local Kummer residue
map is unchanged. Proper base change for $X$ and $D$, followed by the
five lemma in (DC.1), identifies the stalk of $\mathcal M$ with
$H^1(U_s,a)$ and the remaining stalks with $H^q(U_s,a)$. The same
argument commutes with every finite \'etale base change $S'\to S$.
The local compatibility just used can also be checked in a relative
parameter $t$: the residue of a Kummer class $t$ is $1$, both before
and after restriction to the fibre. All schemes in this calculation
are smooth over $k$ and all torsion has order invertible in $k$.

\paragraph{Killing cohomology by $r$-covers of the base.}
If $\mathcal V$ is a finite locally constant $a$-sheaf on $S$ with
$r$-group monodromy image $B$, a class of $H^1(S,\mathcal V)$ is a
continuous crossed homomorphism $c:\pi_1(S)\to V$. The map
\[
 g\longmapsto(c(g),\rho(g))\in V\rtimes B
\]
is a continuous homomorphism to a finite $r$-group. Its kernel defines
a connected finite Galois $r$-cover on which both the sheaf and the
class become trivial.

Also, $H^2(S,a)$ is killed by pullback to any connected \'etale cover
of degree $r$. Such a cover exists because $g(S)\geq1$ and
$H^1(S,\mu_r)\simeq\operatorname{Pic}^0(S)[r]\ne0$.
Under the degree description
$H^2(S,\mu_r)\simeq\operatorname{Pic}(S)/r\simeq\mathbf Z/r$,
pullback multiplies degree by $r$, and is zero. These statements remain
true after every finite \'etale base change, since its genus is still
at least one. A tower of finite Galois $r$-covers is dominated by a
finite Galois $r$-cover of the original base: take a finite normal core,
whose kernel quotient embeds in a product of conjugate $r$-groups.

The first five terms of Leray show that every prescribed class of
$H^1(U_s,a)$ extends to $H^1(U_{S'},a)$ after such a base change.
Indeed, first kill the $r$-group monodromy of $\mathcal M$; the class
then becomes a section of $\mathcal M$. Its sole obstruction to lifting
is its transgression in $H^2(S',a)$, which the preceding degree-$r$
cover kills.

Moreover every class in $H^2(U,a)$ is killed after a finite $r$-cover
of the base. Since $R^2f_*a=0$ and $H^3(S,a)=0$, the Leray filtration
gives the exact sequence
\[
 H^2(S,a)\longrightarrow H^2(U,a)
 \longrightarrow H^1(S,\mathcal M)\longrightarrow0.              \tag{DC.2}
\]
Kill the image in the last term by its affine $r$-group representation
as above. The restricted class then comes from $H^2(S',a)$, and a
further degree-$r$ \'etale cover kills that term. The filtration and
these operations are functorial in the base change.

\paragraph{Central induction for nonabelian torsors.}
We prove the extension assertion by induction on the order of the finite
$r$-group $G$. The trivial group is immediate. Choose a central
subgroup $A\subseteq G$ of order $r$, and put $G_0=G/A$.
Inductively, the quotient of the prescribed fibre $G$-torsor extends to
a $G_0$-torsor on $U_{S_1}$, with a chosen identification on the fibre.
Write $h_0:\pi_1(U_{S_1})\to G_0$ for its continuous homomorphism.
The central extension $1\to A\to G\to G_0\to1$ defines a class
$\omega\in H^2(G_0,A)$, and its pullback
$h_0^*\omega\in H^2_{\mathrm{cont}}(\pi_1(U_{S_1}),A)$ is the
obstruction to a continuous lift to $G$. This assertion follows
directly by choosing a set-theoretic section $G_0\to G$: its factor
set is a two-cocycle, and changing a lift by a one-cochain removes the
defect precisely when that class vanishes. All cochains here are
continuous and have finite values.

For any connected scheme $Z$ used here there is a natural injection
\[
 H^2_{\mathrm{cont}}(\pi_1(Z),A)\hookrightarrow H^2(Z,A).         \tag{DC.3}
\]
To verify it, apply Cartan--Leray to the inverse system of all connected
finite \'etale Galois covers. The first cohomology of that universal
pro-cover is zero: every finite $A$-torsor occurring at a finite level
is trivialized at a higher level. The five-term sequence therefore
gives (DC.3). This argument also shows the natural identification of
continuous and \'etale $H^1$, and is functorial for our base changes.
In particular we have not assumed that the total space is an \'etale
$K(\pi,1)$.

By (DC.2), after a finite $r$-cover of the base the image of
$h_0^*\omega$ in \'etale $H^2$ vanishes. Injectivity in (DC.3) supplies
a continuous lift $h:\pi_1(U_{S_2})\to G$. Its restriction to the
fibre and the prescribed fibre lift have the same image in $G_0$.
Their pointwise ratio, because $A$ is central, is a homomorphism to
$A$, hence a class of $H^1(U_s,A)$. Extend that class after one more
$r$-cover of the base by the preceding $H^1$ argument, and multiply
$h$ by the resulting central character. The adjusted lift restricts
to the prescribed fibre homomorphism. Thus its associated $G$-torsor
is the required extension. If the original fibre torsor is connected,
the total torsor is connected as well, since each connected component
of a finite \'etale cover surjects onto the connected base space and
would otherwise disconnect its fibre.

Finally replace $S'\to S$ by a dominating Galois $r$-cover.
Its pullback gives a normal $r$-index subgroup $H'$ of $H=\pi_1(U)$.
The extended torsor is a finite $r$-quotient of $H'$, with kernel $J$.
The intersection of $\gamma J\gamma^{-1}$ over the finitely many
cosets $H/H'$ is normal in $H$; its quotient is an extension of an
$r$-group by an $r$-group. Restriction to the fibre dominates the
prescribed quotient. Finite quotients separate points in a pro-$r$
group, which proves injectivity.
\end{proof}

\begin{claim}{A.DC2}{Proved}{Right exactness for tame configuration quotients}
In the situation of A.DC1, let $U\to S$ be the open projection.
For every connected finite Galois \'etale cover $W\to U$ tame along
$D$, with group $T$, the image $I\subseteq T$ of $\pi_1(U_s)$ is
normal, and $T/I$ is a quotient of $\pi_1(S)$ through the projection.
Consequently the sequence of tame fundamental groups is exact at its
middle and right terms. The same conclusion holds for any quotient of
$\pi_1(U)$ all of whose finite quotients are tame along $D$.
\end{claim}

\begin{proof}
Normalize $X$ in $W$, obtaining $\bar W\to X$. The normalization is
finite, since $X$ is excellent, and is proper over $S$.
SGA~1, Expos\'e XIII, Proposition 5.2, says that a tame cover of a
regular local scheme with normal-crossings boundary becomes \'etale
after adjoining roots of the boundary parameters of orders equal to
the relevant inertia least common multiples; the orders are invertible.
For our smooth relative divisor with disjoint sections, each strictly
local connected branch is therefore a Kummer chart $z^e=t$, where
$t$ is a relative parameter and $e$ is invertible. This chart is smooth
over $S$, with relative coordinate $z$; its boundary $z=0$ is \'etale
over $S$. Thus $\bar W\to S$ is smooth proper, and
$\bar W-W$ is finite \'etale over $S$.

The Stein factor $S_W$ of $\bar W\to S$ is finite \'etale over $S$.
For clarity, this is the usual proper-smooth connected-component
statement: the geometric fibres are reduced, and the algebra of
global functions on a proper smooth fibre is a product of copies of
its algebraically closed ground field; proper base change and the
Stein factorization identify these algebras in a finite \'etale family.
Equivalently, SGA~1, Expos\'e X, Proposition 1.2 applies to this proper
flat morphism with geometrically separable fibres.
Its fibre over $s$ is the set of connected components of $\bar W_s$.
Deleting finitely many points from a connected smooth curve does not
disconnect it, so this is also the set of components of $W_s$.

The $T$-action extends uniquely to the normalization and to $S_W$.
At $s$, its transitive action on components identifies that finite set
with $T/I$ as a homogeneous set. Since $S_W$ is connected, its base
monodromy is transitive and commutes with $T$. The automorphism group
of the transitive $T$-set $T/I$ is $N_T(I)/I$; it can be transitive
only when $N_T(I)=T$. Hence $I$ is normal. It follows that
$S_W\to S$ is the finite \'etale Galois cover with group $T/I$, and
the quotient map $\pi_1(U)\to T/I$ factors through $\pi_1(S)$.
Surjectivity to the base follows because the pullback of every
connected finite \'etale base cover remains connected. Taking all
finite quotients gives the assertion; images of profinite groups are
closed, so equality of their images in every finite quotient is
equality of the closed subgroups themselves.
\end{proof}

\paragraph{The repaired proof of A.CF12 (paper A.13).}
Let $C/k$ be smooth proper of genus at least two over an algebraically
closed field of characteristic $p>0$. Let $\Sigma$ be arbitrary and
$r\in\Sigma\setminus\{p\}$. Start with a finite quotient $Q$ of
$F_{\mathrm{cr}}$, let $R$ be its normal $r$-subgroup given by the
image of $K_F$, and let $A=Q/R$. The finite quotient $A$ of
$\Delta=\pi_1(C)^\Sigma$ defines a connected finite \'etale proper
Galois cover $C'\to C$. There is no restriction on the $p$-part of
$|A|$. Pull back the two coordinates to obtain
\[
 V=(C'\times C')-\coprod_{a\in A}\Gamma_a.
\]
This is a connected finite \'etale Galois cover of $U$ with group
$A\times A$; its second projection has constant compactification and
the disjoint labelled sections $y\mapsto(ay,y)$. Its fibre is the
cover of $C-P$ corresponding to the kernel of its proper quotient $A$.
The restriction of $Q$ to that fibre has $r$-group image. A.DC1
dominates it by a finite $r$-quotient of $H=\pi_1(V)$, with kernel
$J$. Write $G=\pi_1(U)$ and
$J_0=\bigcap_{g\in G/H}gJg^{-1}$. Since $J$ is normal in $H$,
this finite intersection is normal in $G$, and $H/J_0$ embeds
in the product of the conjugate finite $r$-quotients $H/gJg^{-1}$.
Consequently $T=G/J_0$ fits into
an exact sequence
\[
 1\to R_T\to T\to A\times A\to1,
 \qquad R_T\text{ a finite }r\text{-group}.
\]
Every prime divisor of $|T|$ belongs to $\Sigma$, so $G\to T$
factors through $E=G^\Sigma$. Its map to $A\times A$ factors
through the proper product $\Delta\times\Delta$; hence the image
of $K_E$ lies in $R_T=H/J_0$, and the map factors through
$E_{\mathrm{cr}}$. If a fibre element maps trivially to $T$,
its proper image in $A$ is trivial, so it belongs to the fibre
of $V\to C'$; there it is killed by the chosen quotient of $H$
and therefore by $Q$. Thus the fibre restriction dominates $Q$.
This proves injectivity of
$F_{\mathrm{cr}}\to E_{\mathrm{cr}}$ for every such $\Sigma$.

We spell out the middle exactness, which cannot be inferred from an
unrestricted homotopy sequence for an open morphism in characteristic
$p$. Every inertia group along the diagonal has trivial image in the
proper product. Its image in every finite quotient of
$E_{\mathrm{cr}}$ consequently lies in the image of $K_E$, an
$r$-group. Since $r\ne p$, the corresponding cover is tame along the
diagonal. A.DC2 applies. If $T$ is such a finite quotient and $I$ its
fibre image, $T/I$ comes from the second proper curve. It is a finite
$\Sigma$-group, so this factorization is through $\Delta$.
Thus every element of $\ker(E_{\mathrm{cr}}\to\Delta)$ has image in
$I$ in every finite quotient, and hence belongs to the closed image
of $F_{\mathrm{cr}}$. The reverse inclusion is immediate. Projection
to $\Delta$ is onto, and we obtain
\[
 1\to F_{\mathrm{cr}}\to E_{\mathrm{cr}}
 \overset{\operatorname{pr}_2}{\longrightarrow}\Delta\to1.
\]
This argument proves both the originally asserted relative exactness
and the full finite-domination assertion with quotient $A\times A$.
It uses neither Stix's theorem nor a lift to characteristic zero.
Stix's Proposition 2.7 remains a verified alternative source for the
single-prime relative fibre step, not a dependency of this proof.

For the constant-compactification variant, choose the proper fibre
cover for a prescribed finite fibre quotient, and pull it back along
each of the finitely many section maps. A common proper Galois base
cover trivializing those torsors has group a subgroup of a finite
product of their $\Sigma$-deck groups. On the resulting product the
inverse boundary is disjoint and labelled. A.DC1 and the same finite
core produce an $r$-group by the product of the two proper deck
groups, and A.DC2 gives middle and right exactness. Finally A.DC4
below proves the boundary and arithmetic transport clause of A.CF12.
That proof derives its own geometric fibre input from A.DC1--A.DC2,
so this reference creates no circular dependency.

\begin{claim}{A.DC3}{Proved}{Algebraic survival of the entire $r$-adic cusp inertia}
For $C/k$ smooth proper of genus at least one over an algebraically
closed field of characteristic $p\geq0$, $P\in C(k)$, and $r\ne p$,
the tame $r$-inertia $\mathbf Z_r(1)$ at $P$ injects into
$\pi_1(C-P)^{(r)}$. Consequently it injects into every
$F_{\mathrm{cr}}$ with $r\in\Sigma$.
\end{claim}

\begin{proof}
Choose a nonzero element of $\operatorname{Pic}^0(C)[r]$.
Its Kummer torsor is a connected cyclic \'etale cover $C'\to C$
of degree $r$. There are at least two distinct points $P_1,P_2$ above
$P$. For every $n\geq1$, divisibility of $\operatorname{Pic}^0(C')(k)$
under $r^n$ gives a degree-zero divisor class $[B]$ with
$r^n[B]=[P_1-P_2]$. Choosing a representing divisor, we obtain a
rational function $h$ with
\[
 \operatorname{div}(h)=P_1-P_2-r^nB.
\]
The extension $z^{r^n}=h$ has degree $r^n$ because the valuation of
$h$ at $P_1$ is congruent to $1$ modulo $r^n$. It has inertia of
order $r^n$ at $P_1$ and $P_2$, and is unramified elsewhere outside
the preimage of $P$: every other valuation is divisible by $r^n$,
and the residue fields are algebraically closed. Taking the compositum
of its conjugates under the cyclic deck group of $C'\to C$ produces
a Galois extension of $k(C)$ with $r$-group Galois group, ramified
only above $P$, whose inertia surjects onto a cyclic group of order
$r^n$.

These are finite \'etale $r$-covers of $C-P$ for every $n$.
The image of the local tame $\mathbf Z_r(1)$ therefore has unbounded
$r$-power finite quotients; its kernel, a closed subgroup of
$\mathbf Z_r$, must be zero. All these quotients factor through
$F_{\mathrm{cr}}$, since their images of $K_F$ are $r$-groups.
The Kummer and Picard facts used here are algebraic; see
\citet[Propositions 14.1--14.2]{CF_Milne2013}. In particular the proof
does not use the topological presentation of a punctured surface.
\end{proof}

\begin{claim}{A.DC4}{Proved}{Diagonal boundary and arithmetic transport, with $p$ allowed in $\Sigma$}
Let $C/k$ be smooth proper geometrically connected of genus at least
two over a finite field of characteristic $p$, and let
$r\in\Sigma\setminus\{p\}$. Let $\Pi_C$ be the geometrically
pro-$\Sigma$ arithmetic fundamental group, retaining the whole $G_k$.
Let $\mathcal E_{\mathrm{cr}}$ be the arithmetic relative
cuspidally pro-$r$ configuration group. There is a diagonal boundary
subgroup $\mathcal D\subseteq\mathcal E_{\mathrm{cr}}$ with exact
sequence
\[
 1\to\mathbf Z_r(1)\to\mathcal D\to\Pi_C\to1.
\]
Its map to $\Pi_C\times_{G_k}\Pi_C$ is diagonal.
For each $P\in C(k)$, with compatible paths, the arithmetic group
$\Pi_{C-P,\mathrm{cr}}$ equals
$\operatorname{pr}_2^{-1}(s_P(G_k))$, and its actual cusp subgroup
is $\mathcal D\cap\operatorname{pr}_2^{-1}(s_P(G_k))$.
If the proper sections at $P$ and $Q$ are geometrically conjugate,
their punctured arithmetic groups are isomorphic over the identity of
$\Pi_C$ and $G_k$, with the first actual cusp taken to a geometric
conjugate of the second.
\end{claim}

\begin{proof}
We first obtain the geometric fibre input directly from A.DC1--A.DC2.
For a finite fibre quotient $Q$, divide by its cuspidal $r$-subgroup
to obtain the proper quotient $A$, and use the $A\times A$ deck-graph
cover $V$. A.DC1 gives a finite $r$-quotient of $\pi_1(V)$ detecting
the restriction of $Q$. The intersection of its kernel's conjugates
over $A\times A$ has quotient in $\pi_1(V)$ embedded in a finite
product of $r$-groups. Its total quotient is an $r$-group by
$A\times A$, factors through $E_{\mathrm{cr}}$, and dominates $Q$
on the fibre: an element killed in the total quotient first has
trivial proper $A$-image and is then killed by the chosen fibre
quotient. Thus the fibre map is injective. Its finite total quotients
are tame along the diagonal, since diagonal inertia has $r$-power
image; A.DC2 gives middle and right exactness. This establishes
the required geometric short exact sequence within the present
proof, independently of A.CF12.

Put $X=C\times_k C$, $D=\operatorname{diag}(C)$, and
$N=\mathcal O_X(D)|_D$. For $n=r^a$, the root stack
$\mathfrak X_n=\sqrt[n]{(X,D)}$ is defined locally, if $D=(t)$, by
\[
 \bigl[\operatorname{Spec}(R[z]/(z^n-t))/\mu_n\bigr].
\]
Its boundary $\mathfrak G_n$ is the gerbe of $n$-th roots of $N$ on
$D$. We use this language for descent of the Kummer charts; we do not
attribute a statement about stacks to SGA~1.
Specifically, SGA~1 XIII, Proposition 5.2, extends any finite tame
cover whose inertia orders divide $n$ to an \'etale cover of each
Kummer chart. Uniqueness of normalization extends the descent
isomorphisms, including their $\mu_n$-equivariance, from the dense open
complement. They glue to a finite \'etale cover of $\mathfrak X_n$.
Conversely, restriction of a finite \'etale cover of that stack to
$X-D$ has inertia of order dividing $n$. This proves precisely the
cover comparison used below, by finite \'etale descent.

A finite \'etale cover of $\mathfrak G_n$ on which its stabilizer
$\mu_n$ acts trivially descends to $D$: this is ordinary descent on
local trivializations $B\mu_n\times D_i$. Pullbacks of covers of $D$
are fully faithful. The corresponding fundamental group sequence is
therefore right exact:
\[
 \mu_n\longrightarrow\pi_1(\mathfrak G_n)
 \longrightarrow\pi_1(C)\longrightarrow1.                       \tag{DC.4}
\]
No injectivity of the finite stabilizer is assumed. The transition
maps for $r^{a+1}\to r^a$ are onto on the base and on the stabilizer
images, hence onto on the middle groups. Compact inverse limits of
(DC.4) give
\[
 \mathbf Z_r(1)\longrightarrow\mathcal T
 \longrightarrow\pi_1(C)\longrightarrow1.                      \tag{DC.5}
\]
Every finite quotient of $\mathcal E_{\mathrm{cr}}$ is tame with
$r$-power inertia, for the reason just given.
Restriction of its extended root-stack cover to the boundary produces
compatible maps from $\mathcal T$; thus there is a continuous map
$\mathcal T\to\mathcal E_{\mathrm{cr}}$ over the diagonal proper
map. Define $\mathcal D$ as its closed image.

Here is the precise quotient argument at the proper geometric base.
Write $\mathcal T_{\bar k}$ for the geometric kernel of $\mathcal T$,
and replace it by its maximal pro-$\Sigma$ quotient, leaving $G_k$
unchanged. The resulting arithmetic quotient $\mathcal T_\Sigma$
has a right exact sequence
\[
 \mathbf Z_r(1)\longrightarrow\mathcal T_\Sigma
 \longrightarrow\Pi_C\longrightarrow1.                        \tag{DC.6}
\]
Indeed, on geometric groups (DC.5) is an extension of $\pi_1(C_{\bar k})$
by the inertia image. After maximal pro-$\Sigma$ quotient, quotienting
out this image represents exactly the finite $\Sigma$-quotients of
$\pi_1(C_{\bar k})$; this is right exactness of completion, not an
assumption that it preserves arbitrary kernels. The map to
$\mathcal E_{\mathrm{cr}}$ factors through $\mathcal T_\Sigma$
because its geometric image is pro-$\Sigma$. Consequently
$\mathcal D\to\Pi_C$ is onto, with kernel exactly its inertia image,
a quotient of $\mathbf Z_r(1)$.

The fibre immersion $C-P\to X-D$, $x\mapsto(x,P)$, extends to a
map of pairs transverse to $D$ with multiplicity one. In the local
Kummer charts it takes the diagonal inertia to the actual cusp inertia
with that multiplicity. A.DC3 shows that the latter is the whole
$\mathbf Z_r(1)$ in $F_{\mathrm{cr}}$, and the geometric injection
just proved embeds this fibre group into $E_{\mathrm{cr}}$. Thus inertia in (DC.6) has
zero kernel in $\mathcal D$, proving the displayed short exact
sequence of the statement. This argument also works for $r=2$.

The actual arithmetic fibre induces a homomorphism
\[
 \Pi_{C-P,\mathrm{cr}}\longrightarrow
 \mathcal E_{\mathrm{cr}}\times_{\Pi_C,s_P}G_k.
\]
It is an isomorphism on geometric kernels by the exact sequence
proved above and the identity
on $G_k$; comparison of short exact sequences makes it an isomorphism.
Since $s_P$ is injective, the target is
$\operatorname{pr}_2^{-1}(s_P(G_k))$. The local punctured henselian
trait at $P$ gives the actual cusp decomposition subgroup. Its
inertia is $\mathbf Z_r(1)$ by A.DC3 and its residue quotient is the
whole $G_k$, retained injectively by the proper point section.
Its map into $\mathcal D$ is the identity on these inertia and
residue groups. The short exact sequence comparison therefore
identifies it with the full inverse image of $s_P(G_k)$ in
$\mathcal D$. Paths may be chosen in the same connected boundary
tower for all points under comparison, so this statement uses one
common $\mathcal D$ and does not assume point separation.

Finally suppose $\gamma\in\Delta$ satisfies
$\gamma s_P(g)\gamma^{-1}=s_Q(g)$ for every $g\in G_k$.
Choose a geometric element $d\in\mathcal D$ mapping to $\gamma$.
Conjugation by $d$ identifies the two fibre subgroups and their actual
cusps. On the first proper coordinate it induces
$\operatorname{Inn}(\gamma)$. The geometric subgroup of the $Q$-fibre
surjects onto $\Delta$ under that first coordinate, so choose in it
$e$ mapping to $\gamma^{-1}$. Then conjugation by $ed$ is over the
identity of $\Pi_C$ and $G_k$. It carries the first actual cusp to
$e$ times the second cusp times $e^{-1}$, a geometric conjugate.
All choices change the resulting identifications only by the geometric
inner conjugacies explicitly allowed in the statement.
\end{proof}

\paragraph{Dependency conclusion for G1(ii).}
The geometric proof of A.CF12 is A.DC1 plus the finite proper-cover
and finite-core construction written out above, followed by A.DC2
for exactness at the middle term. Its boundary and arithmetic clause
uses A.DC4. The latter independently derives its geometric input
from A.DC1--A.DC2 and uses A.DC3 for inertia; it does not depend on
A.CF12. These four arguments are in characteristic $p$,
including when $p\in\Sigma$; they invoke no lifting theorem and none
of SGA~1 III 7.3--7.4 or XIII 2.9--2.11. SGA~1 XIII 5.2 is used for
local tame Kummer charts only. In contrast, the original Stix route
uses XIII 4.3--4.4, and the latter explicitly cites XIII 2.9 for its
cohomological hypotheses. It proves the desired geometry without an
actual lift, but does not by itself satisfy the stricter literal
source-exclusion condition. The present direct central-induction proof
removes that dependency rather than concealing it.
